% ===========================================================================
% 00_frontmatter.tex - preface, scope, guidelines, report format, metrics
% ===========================================================================

\chapter*{Preface}
\phantomsection
\addcontentsline{toc}{chapter}{Preface}
This manual is designed around the algorithm list supplied for the Machine Learning
assignment. It turns that list into a structured laboratory sequence emphasizing
implementation, experimentation, model evaluation and interpretation rather than only
code execution. Students should run the programs, record outputs, vary important
hyperparameters, and explain why the observed results change.

This \textbf{completed edition} goes one step further: every experiment is implemented
and executed on real, public datasets, with the observed results, discussion and the
complete runnable code included, so each laboratory can be reproduced standalone.

\chapter*{Important Terminology Correction}
\phantomsection
\addcontentsline{toc}{chapter}{Important Terminology Correction}
The supplied list contains ``Random Forest Regression (RFR)'' and ``Random Forest
Regression (RFC)''. In standard machine-learning terminology, RFC normally means
Random Forest Classifier. This manual therefore treats \textbf{RFR} as Random Forest
Regressor and \textbf{RFC} as Random Forest Classifier.

\chapter*{Scope of the Laboratory}
\phantomsection
\addcontentsline{toc}{chapter}{Scope of the Laboratory}
\begin{longtable}{|p{0.28\textwidth}|p{0.66\textwidth}|}
\hline
\rowcolor{mlblue!8}\textbf{Module} & \textbf{Algorithms / Topics} \\
\hline
\endhead
Clustering & K-means, Modified K-means, Hierarchical Clustering, Fuzzy C-means \\
\hline
Density-based learning & DBSCAN, HDBSCAN \\
\hline
Semi-supervised learning & Self-training \\
\hline
Ensemble learning & Random Forest Regression, Random Forest Classification, XGBoost, AdaBoost, CatBoost \\
\hline
Neural networks & MLP, RNN (LSTM) \\
\hline
Unsupervised neural learning & SOM \\
\hline
Probabilistic sequence model & HMM \\
\hline
Margin-based learning & SVM (SVC and SVR) \\
\hline
Generative / foundation models & LLM \\
\hline
Non-parametric regression & GRNN \\
\hline
\end{longtable}

\chapter*{General Learning Outcomes}
\phantomsection
\addcontentsline{toc}{chapter}{General Learning Outcomes}
\begin{itemize}
  \item Implement and execute major machine-learning algorithms using Python.
  \item Select an appropriate algorithm based on data type, task and assumptions.
  \item Preprocess data, handle categorical variables, scale numerical features and split datasets correctly.
  \item Evaluate clustering with internal measures and supervised models with appropriate metrics.
  \item Interpret model behavior through decision boundaries, feature importance, learning curves and error analysis.
  \item Compare algorithms using a fair experimental protocol rather than relying on a single score.
  \item Document reproducible experiments, hyperparameters, random seeds, observations and limitations.
\end{itemize}

\chapter*{Recommended Software Environment}
\phantomsection
\addcontentsline{toc}{chapter}{Recommended Software Environment}
\begin{codebox}[title={Core and specialized libraries used by the manual},listing options={language=bash}]
# Core
pip install numpy pandas matplotlib seaborn scikit-learn
# Specialized algorithms
pip install scikit-fuzzy hdbscan hmmlearn minisom
# Gradient boosting libraries
pip install xgboost catboost
# LLM / transformer experiment
pip install transformers torch
\end{codebox}

\chapter*{Laboratory Guidelines}
\phantomsection
\addcontentsline{toc}{chapter}{Laboratory Guidelines}
\begin{longtable}{|p{0.30\textwidth}|p{0.34\textwidth}|p{0.30\textwidth}|}
\hline
\rowcolor{mlblue!8}\textbf{Before Lab} & \textbf{During Lab} & \textbf{After Lab} \\
\hline
\endhead
Read theory and algorithm steps. & Run the baseline program. & Submit code and results. \\
\hline
Know the dataset and target. & Record hyperparameters and seed. & Explain observations. \\
\hline
Predict expected behavior. & Perform at least one controlled variation. & Answer viva questions. \\
\hline
Prepare key formulas. & Check train/test leakage. & State limitations and improvements. \\
\hline
\end{longtable}

\chapter*{Minimum Report Format}
\phantomsection
\addcontentsline{toc}{chapter}{Minimum Report Format}
\begin{enumerate}
  \item Experiment title and student information.
  \item Problem statement and objective.
  \item Dataset description and preprocessing.
  \item Algorithm principle and pseudocode / flowchart.
  \item Python implementation.
  \item Hyperparameters and experimental setup.
  \item Results: tables, plots and relevant metrics.
  \item Discussion: what the results mean and why they occurred.
  \item Conclusion and limitations.
  \item Answers to viva questions.
\end{enumerate}

\chapter*{Recommended Evaluation Metrics}
\phantomsection
\addcontentsline{toc}{chapter}{Recommended Evaluation Metrics}
\begin{longtable}{|p{0.26\textwidth}|p{0.68\textwidth}|}
\hline
\rowcolor{mlblue!8}\textbf{Task} & \textbf{Suggested Metrics} \\
\hline
\endhead
Classification & Accuracy, Precision, Recall, F1-score, ROC-AUC, confusion matrix \\
\hline
Regression & MAE, MSE, RMSE, R\textsuperscript{2} \\
\hline
Clustering & Silhouette score, Calinski--Harabasz, Davies--Bouldin; visual inspection \\
\hline
Sequence modeling & Negative log-likelihood / log-likelihood, sequence accuracy where appropriate \\
\hline
Language generation & Task-specific automatic metrics plus qualitative / error analysis \\
\hline
\end{longtable}
